\section{İlgili Çalışmalar}
\label{sec:related}

\subsection{HiperKvasir ve gastrointestinal endoskopi sınıflandırması}

HiperKvasir, gastrointestinal endoskopi görüntüleri ve videoları için yayımlanmış
geniş kapsamlı bir açık veri kümesidir~\cite{borgli2020hyperkvasir}. Bu raporda
kullanılan denetimli alt küme, 23 sınıfa ayrılmış 10.662 etiketli görüntüden oluşur.
Sınıflar anatomik işaretleri, patolojik bulguları, tedaviyle ilişkili görüntüleri ve
görüntü kalitesi kategorilerini birlikte içerdiği için problem sadece patoloji-normal
ayrımına indirgenemez. Ayrıca sınıf dağılımı belirgin biçimde dengesizdir; çok düşük
destekli sınıflar, genel doğruluğu yüksek görünen modellerde bile makro-F1'i sınırlayan
ana etkenlerden biridir.

HiperKvasir ve yakın gastrointestinal endoskopi veri kümeleri üzerinde çeşitli derin
öğrenme yaklaşımları raporlanmıştır. HSW-ViT, yerel ve uzun menzilli bağımlılıkları
birlikte ele almak için Swin tabanlı hibrit kaydırılmış pencere tasarımını
incelemiştir~\cite{wang2023hswvit}. GastroViT, 23 sınıflı HiperKvasir bağlamında
hafif Vision Transformer tabanlı bir topluluk ve açıklanabilirlik analizleri
sunmuştur~\cite{gastrovit2025}. EffiMix ise EfficientNet ve özel evrişimsel
öznitelikleri birleştiren bir füzyon yaklaşımı olarak aynı veri ailesinde yüksek
toplu başarılar bildirmiştir~\cite{ramamurthy2022effimix}. Bu çalışmalar, endoskopi
sınıflandırmasında hem yerel doku ipuçlarının hem de daha geniş bağlamsal temsillerin
önemli olduğunu gösterir. Ancak kullanılan sınıf alt kümeleri, bölme protokolleri,
ölçütler ve eğitim düzenleri farklıdır; bu nedenle bu rapordaki karşılaştırmalar
doğrudan sıralama amacı taşımaz, literatür sonuçları bağlamsal referans olarak ele
alınır.

\subsection{Görüntü dönüştürücüleri ve tıbbi görüntü sınıflandırması}

Vision Transformer (ViT), görüntüyü sabit boyutlu yamalara ayırıp bu yamaları standart
bir dönüştürücü kodlayıcıya girdi olarak verir; bu tasarım, evrişimsel ağlara göre daha
zayıf görüntüye özgü önyargılarla küresel patch-token ilişkilerini modellemeye
dayanır~\cite{dosovitskiy2021vit}. Swin Transformer, öz-dikkati tüm görüntü yerine
kaydırılmış yerel pencerelerde hesaplayarak hiyerarşik ve çok ölçekli temsiller üretir;
bu yapı özellikle farklı ölçeklerde görülebilen endoskopik bulgular için yararlı bir
tamamlayıcı bakış sağlar~\cite{liu2021swin}. BEiT ise maskeli görüntü modellemesiyle
ön-eğitim görür ve maskelenen görsel token'ları tahmin etmeye dayalı öz-denetimli
ön-eğitim çizgisini takip eder~\cite{bao2022beit}. Tıbbi görüntülemede etiketleme
maliyetli olduğundan, böyle ön-eğitim rejimleri sınırlı etiketli veriyle aktarım
öğrenmesini daha anlamlı hale getirir.

Bu raporda ViT-B/16, Swin-T ve BEiT-B/16 omurgalarının birlikte seçilmesi bu
tamamlayıcılık fikrine dayanır: ViT-B/16 daha küresel patch-token temsili sağlar,
Swin-T hiyerarşik kaydırılmış pencere yapısıyla yerel-çok ölçekli bağlamı yakalar,
BEiT-B/16 ise maskeli görüntü modellemesiyle edinilmiş ön-eğitim bilgisini taşır. Üç
omurga da önceden eğitilmiş ağırlıklarla kullanılır ve her biri 768 boyutlu
havuzlanmış özellik üretir; bu özellikler 512 boyutlu dal projeksiyonuna indirgenip
füzyon katmanına aktarılır. Literatürdeki ViT ailesi çalışmalarından farklı olarak bu
rapor, tek bir mimarinin üstünlüğünü iddia etmek yerine aynı resmi 5-katlı protokol
altında üç farklı dönüştürücü önyargısının tekli, ikili ve üçlü katkısını ablation
mantığıyla değerlendirir.

\subsection{Özellik füzyonu ve çoklu omurga yaklaşımları}

Çoklu omurga veya çok kaynaklı özellik füzyonu, farklı ağların öğrendiği temsilleri
karar katmanından önce birleştirerek tek bir temsil uzayı oluşturmayı amaçlar. GI
endoskopi bağlamında CNN özniteliklerinin hibrit bir füzyon hattında birleştirilmesi
Ahmed ve arkadaşları tarafından incelenmiştir~\cite{ahmed2023hybridfusion}. HEMF,
tıbbi görüntü sınıflandırmasında yerel ve küresel çok ölçekli özellikleri hiyerarşik
dikkat mekanizmalarıyla birleştiren daha genel bir tıbbi görüntü füzyon örneği
sunar~\cite{he2025hemf}. CNN ve Swin temsillerinin birlikte kullanıldığı
endoskopi/video-kapsül çalışmaları da, evrişimsel yerel desenler ile dönüştürücü
bağlamsal temsillerinin tamamlayıcı olabileceğini göstermektedir~\cite{subedi2024cnnswin}.
Kaynak kısıtlı dağıtım senaryolarında ise hafif ViT ailesi modelleri, performans ile
model boyutu arasındaki dengeyi ayrıca gündeme getirir~\cite{varam2024edgevit}.

Bu rapordaki füzyon tasarımı daha kontrollü tutulmuştur. Her omurga aynı 512 boyutlu
projeksiyon biçiminden geçirilir; sınıflandırıcı aşamasında ayrı bir klasik
sınıflandırıcı ailesine geçilmez, bütün deneylerde sabit MLP sınıflandırıcı korunur.
Böylece performans farkları sınıflandırıcı ailesinden çok omurga seçimi, aktarım modu
ve füzyon biçimine bağlanabilir. Birleştirme
(concatenation) ve öğrenilebilir ağırlıklı füzyon pratik ana seçeneklerdir. GMU,
farklı modalite veya dalların katkısını öğrenilen kapılarla ayarlayan ilkeli bir
füzyon fikri sunar~\cite{arevalo2017gmu}; bu raporda ise GMU ana yöntem değil, ek
ablation olarak değerlendirilmiş ve üçlü ağırlıklı füzyonu geçemediği için negatif
sonuç olarak raporlanmıştır.

\subsection{Dengesizlik, değerlendirme ve yorumlanabilirlik}

HiperKvasir'in 23 sınıflı alt kümesindeki sınıf dengesizliği, değerlendirme ölçütü
seçimini doğrudan etkiler. Doğruluk, yüksek destekli sınıflarda iyi performans gösteren
bir model için iyimser bir tablo verebilir; makro-F1 ise her sınıfı eşit ağırlıkla
ortaladığı için nadir sınıflardaki hataları görünür kılar. Bu nedenle raporda ana
ölçüt makro-F1, doğruluk ise destekleyici ölçüt olarak kullanılır. Aynı gerekçeyle
resmi 5-katlı ve sızıntısız OOF değerlendirme protokolü, tek fold sonuçlarının veya
aynı görüntünün farklı aşamalarda karışmasının doğurabileceği iyimserliği azaltmak
için merkezîdir. Model karşılaştırmalarında McNemar testi, aynı örnekler üzerindeki
eşli doğru/yanlış kararlarını sınar; bu test doğruluk düzeyinde yorumlanır ve makro-F1
için eşdeğer bir anlamlılık testi olarak sunulmaz~\cite{dietterich1998mcnemar}.

Dengesiz veri için literatürde zor örneklere odaklanan focal loss~\cite{lin2017focal}
ve sınıf öncüllerini logit düzeyinde düzelten yöntemler~\cite{menon2020logitadjust}
gibi seçenekler bulunur. Bu raporda sınıf dengesizliği öncelikle eğitim örneklemesi,
makro-F1 merkezli değerlendirme ve sınıf bazlı hata analiziyle ele alınmıştır. Logit
düzeltmesi ayrıca denenmiş, fakat mevcut eğitim düzeni altında makro-F1'i düşürdüğü
için final model seçimini değiştirmemiştir. Bu sonuç, dengesizlik yöntemlerinin veri
protokolü ve eğitim rejimiyle birlikte değerlendirilmesi gerektiğini gösterir.

Yorumlanabilirlik tarafında attention rollout, dönüştürücü katmanlarındaki dikkat
matrislerini kalıntı bağlantılarla birlikte izleyerek token'lar arası bilgi akışını
yaklaşıklar~\cite{abnar2020rollout}. Grad-CAM++ sınıfa özgü gradyan ağırlıklarıyla
ısı haritaları üretir ve Grad-CAM'in yerelleştirme fikrini genişletir
~\cite{selvaraju2017gradcam,chattopadhyay2018gradcampp}. Rapordaki attention rollout,
Grad-CAM++ ve UMAP görselleştirmeleri, modelin hangi görüntü bölgelerine ve hangi
özellik uzayı ayrışmalarına duyarlı olduğuna dair nitel destek sağlar. Bu görseller
klinik doğrulama veya nicel açıklanabilirlik kanıtı olarak değil, hata analizi ve
temsil davranışını yorumlama aracı olarak kullanılmalıdır.

Bu literatür çerçevesinde raporun katkısı, farklı protokollerde bildirilmiş tekil
literatür sayılarıyla yarışmak değil; aynı resmi 5-katlı OOF protokolü altında üç
ViT ailesi omurgasını, tekli/ikili/üçlü ablation düzenini, dondurulmuş özellik çıkarımı
ile son blok ince ayarını, MLP-only sınıflandırıcıyı ve makro-F1 merkezli raporlamayı
bir arada kontrol etmektir. Final model seçimi de bu kontrollü kurgu içinde,
ViT-B/16 + Swin-T + BEiT-B/16 ağırlıklı füzyonunun 5-fold ve pooled OOF sonuçlarına
dayanır.
